
\chapter{Load Balancers and Anycast}
Load balancing presents an interesting design challenge. The canonical mental model of a load balancer is a single entry point that forwards incoming requests to one of several backend services, with minimal added latency. Since it is the single point of entry, it must handle all incoming traffic. For instance, Google uses the IP address 8.8.8.8, yet the entire world can contact it without overwhelming a single machine. Understanding how this is achieved requires revisiting the assumption that an IP address maps to a single physical server.

The answer is that the assumption of a one-to-one mapping between an IP address and a physical server is incorrect. In practice, a load balancer is a cluster of machines that share the same IP address, with each machine accessible through a different router.

This technique is called \textit{Anycast} routing. Each machine in the cluster advertises the same IP address to the network via BGP (Border Gateway Protocol). When a client sends a packet to that IP, the routing infrastructure delivers it to whichever cluster node is topologically closest to the client. Different clients around the world therefore reach different physical machines, even though they all contact the same IP address. This distributes the load geographically without requiring the client to be aware of the cluster at all.

A simpler alternative is \textit{DNS-based} load balancing, where the DNS server returns a different IP address for each query, cycling through the addresses of the backend servers. This approach does not require shared IPs, but it is less flexible: clients and resolvers cache DNS responses, so a client may continue contacting a machine that has gone offline until the cached record expires.

\section{BGP and Anycast in Detail}

Anycast is not a dedicated feature of BGP. It is an emergent consequence of running BGP normally while multiple machines simultaneously advertise the same IP prefix. BGP was designed for Unicast, where each prefix has a single origin, but nothing in the protocol prevents multiple origins from advertising the same block. Each BGP router receives all of these advertisements and independently selects the single best path according to its routing policy, considering factors such as AS path length and local preference. Crucially, routers in different parts of the network have different topological perspectives, so they select different best paths. A router in Tokyo will prefer the Tokyo Anycast node; a router in Frankfurt will prefer the Frankfurt one. From each router's local view there is one route, but globally the same prefix resolves to different physical machines depending on where the query originates.

\paragraph{Session maintenance and failover.}
BGP sessions are kept alive with periodic \textit{KEEPALIVE} messages, sent by default every 30 seconds. Each session carries a \textit{hold timer}, defaulting to 90 seconds. If no KEEPALIVE or UPDATE message arrives within the hold timer window, the router declares the session dead and withdraws all routes learned from that peer. When an Anycast node goes offline, it stops sending KEEPALIVEs; after the hold timer expires, its peer withdraws the route and propagates a BGP UPDATE message through the network. Other routers then reconverge to the next-best path toward the prefix.

This reconvergence is not instantaneous. Depending on network size and policy, failover may take anywhere from tens of seconds to a few minutes. This delay has an important implication for stateful protocols such as TCP: if a client is mid-connection when BGP reconverges to a different Anycast node, the connection breaks, because the new node holds no session state for that client. For this reason, Anycast is best suited to stateless protocols. DNS is the canonical example: each query is independent, so routing a query to a different node on failover is transparent to the client. Google's 8.8.8.8 DNS resolver is deployed precisely this way.

\paragraph{Multiple advertisements for the same prefix.}
A BGP router can receive advertisements for the same prefix from multiple peers simultaneously, each representing a path to a different Anycast node. The router stores all of these in its \textit{Routing Information Base} (RIB), but installs only the single best path in the forwarding table. The best path is selected through a deterministic tie-breaking process, evaluating candidates in order: local preference, AS path length, origin type, MED (Multi-Exit Discriminator), and finally the router ID as a tiebreaker of last resort. The non-selected paths are not discarded; they remain in the RIB as backup routes. If the best path is withdrawn, the router immediately promotes the next-best path from the RIB without waiting for a new advertisement, making failover faster than if no alternatives were cached.

\paragraph{IP hijacking.}
The same mechanism that enables Anycast also creates a significant security vulnerability. Because BGP has no built-in authentication, any Autonomous System (AS) can advertise any prefix, including one it does not own. If a malicious AS advertises a more specific prefix (i.e., a longer subnet mask) than the legitimate owner, BGP's preference for specificity means that routers near the attacker will prefer the malicious route. Traffic destined for the victim is then redirected to the attacker, a technique known as a \emph{BGP hijack}.

A notable example occurred in 2010, when China Telecom briefly advertised routes for large portions of the internet, redirecting traffic through Chinese routers for approximately 18 minutes. In 2018, traffic destined for Amazon Route 53 DNS servers was hijacked using this technique to steal cryptocurrency.

The standard mitigation is \emph{Resource Public Key Infrastructure} (RPKI), which allows IP prefix owners to cryptographically sign route origin authorizations (ROAs), stating which AS is permitted to advertise a given prefix. Routers that enforce RPKI can reject advertisements that lack a valid ROA. Adoption has grown substantially but remains incomplete, meaning BGP hijacking is still a practical threat.

% \section{Gateway Types}

% Load balancers are commonly classified by the OSI layer at which they operate. The two principal types are Layer 4 and Layer 7 load balancers, each with distinct capabilities and trade-offs.

% \paragraph{Layer 4 (Transport Layer).}
% A Layer 4 load balancer operates on IP addresses and TCP/UDP port numbers, without inspecting the contents of packets. When a connection arrives, the balancer selects a backend server and rewrites the destination IP and port in the packet headers, then forwards it. Because it never reads above the transport layer, it does not need to terminate the TCP connection itself: the connection runs directly between the client and the backend. This makes Layer 4 balancers very fast and lightweight, as the forwarding decision is made once per connection rather than once per request. The trade-off is limited routing intelligence: all packets in a flow go to the same backend, and the balancer cannot distinguish between, for example, a request for a static image and an API call within the same TCP session.

% For instance, AWS uses a tuple of 5 values: (client ip, client port, server ip, server port, and protocol). It applies a hash on the tuple to determine which underlying server should handle the incoming traffic. 
% Usually, it has specialized hardware, and the servers in the same VPC allows for TCP/UDP connection handoff.

% \paragraph{Layer 7 (Application Layer).}
% A Layer 7 load balancer operates at the application layer and can inspect the full content of requests, including HTTP headers, URLs, cookies, and request bodies. It terminates the client's TCP connection, reads the request, makes a routing decision based on its content, and opens a new connection to the selected backend. This enables much richer routing policies: requests to \texttt{/api/} can be forwarded to one pool of servers while requests to \texttt{/static/} go to another; a cookie can pin a user to a specific backend for session affinity; traffic can be shifted between versions of a service for canary deployments. The cost is higher per-request overhead, since the balancer must fully parse each request and maintain two separate TCP connections per client.